\documentclass[../../../ps_q.tex]{subfiles}

\begin{document}
    \setcounter{figure}{0}
    図1のように，真空中に，極板面積$S$，極板間隔$d$の平行平板コンデンサがあり，
    電荷$Q$を与えてある．真空の誘電率を$\varepsilon_0$とし，端の効果は無視する．

    \begin{enumerate}
        \item 電位差（電圧）$V_1$，および容量$C_1$を求めよ．
        \item 図\ref{fig:fig_em_cx_06_01_q}のコンデンサの極板間を
            比誘電率$\varepsilon_r$の誘電体で満たす（図\ref{fig:fig_em_cx_06_02_q}）．このときの電位差
            （電圧）$V_2$，および容量$C_2$を求めよ．
        \item 図\ref{fig:fig_em_cx_06_01_q}のコンデンサの極板間のちょうど下半分の領域を
            比誘電率$\varepsilon_r$の誘電体で満たす（図\ref{fig:fig_em_cx_06_03_q}）．このときの電位差
            （電圧）$V_3$，および容量$C_3$を求めよ．
        \item 図\ref{fig:fig_em_cx_06_01_q}のコンデンサの極板間のちょうど右半分の領域を
            比誘電率$\varepsilon_r$の誘電体で満たす（図\ref{fig:fig_em_cx_06_04_q}）．このときの電位差
            （電圧）$V_4$，および容量$C_4$を求めよ．
    \end{enumerate}
    \vspace{10pt}
    \begin{minipage}{\linewidth}
        \setlength{\abovecaptionskip}{2pt}
        \centering
            \includegraphics{\subfix{fig_em_cx_06/fig_em_cx_06_01_q.pdf}}
            \captionof{figure}{} \label{fig:fig_em_cx_06_01_q}
            \vspace{12pt}
            \includegraphics{\subfix{fig_em_cx_06/fig_em_cx_06_02_q.pdf}}
            \captionof{figure}{} \label{fig:fig_em_cx_06_02_q}
            \vspace{12pt}
            \includegraphics{\subfix{fig_em_cx_06/fig_em_cx_06_03_q.pdf}}
            \captionof{figure}{} \label{fig:fig_em_cx_06_03_q}
            \vspace{12pt}
            \includegraphics{\subfix{fig_em_cx_06/fig_em_cx_06_04_q.pdf}}
            \captionof{figure}{} \label{fig:fig_em_cx_06_04_q}
    \end{minipage}
\end{document}






